%!TeX root=../pridetop.tex
\chapter[Chapter \thechapter]{}
	
	\begin{figure}[t!]
\centering
\includegraphics[width=\linewidth]{13top}
\captionlistentry{Headpiece to Chapter \thechapter}
\end{figure}

\lettrine[ante=`,lines=4,image=true,findent=5pt,realheight=true]{initials/chap13i}{} hope, my dear,' said Mr Bennet to his wife, as they were at breakfast the next morning, »that you have ordered a good dinner to-day, because I have reason to expect an addition to our family party.«

\zz
»Who do you mean, my dear? I know of nobody that is coming, I am sure, unless Charlotte Lucas should happen to call in; and I hope \textit{my} dinners are good enough for her. I do not believe she often sees such at home.«

»The person of whom I speak is a gentleman and a stranger.«

Mrs Bennet's eyes sparkled. »A gentleman and a stranger! It is Mr Bingley, I am sure. Why, Jane—you never dropped a word of this—you sly thing! Well, I am sure I shall be extremely glad to see Mr Bingley. But—good Lord! how unlucky! there is not a bit of fish to be got to-day. Lydia, my love, ring the bell. I must speak to Hill this moment.«

»It is \textit{not} Mr Bingley,« said her husband; »it is a person whom I never saw in the whole course of my life.«

This roused a general astonishment; and he had the pleasure of being eagerly questioned by his wife and five daughters at once.

After amusing himself some time with their curiosity, he thus explained:— »About a month ago I received this letter, and about a fortnight ago I answered it; for I thought it a case of some delicacy, and requiring early attention. It is from my cousin, Mr Collins, who, when I am dead, may turn you all out of this house as soon as he pleases.«

»Oh, my dear,« cried his wife, »I cannot bear to hear that mentioned. Pray do not talk of that odious man. I do think it is the hardest thing in the world, that your estate should be entailed away from your own children; and I am sure, if I had been you, I should have tried long ago to do something or other about it.«

Jane and Elizabeth attempted to explain to her the nature of an entail. They had often attempted it before: but it was a subject on which Mrs Bennet was beyond the reach of reason; and she continued to rail bitterly against the cruelty of settling an estate away from a family of five daughters, in favour of a man whom nobody cared anything about.

»It certainly is a most iniquitous affair,« said Mr Bennet; »and nothing can clear Mr Collins from the guilt of inheriting Longbourn. But if you will listen to his letter, you may, perhaps, be a little softened by his manner of expressing himself.«

»No, that I am sure I shall not: and I think it was very impertinent of him to write to you at all, and very hypocritical. I hate such false friends. Why could not he keep on quarrelling with you, as his father did before him?«

»Why, indeed, he does seem to have had some filial scruples on that head, as you will hear.«

\begin{quotation}
\begin{flushright}
Hunsford, near Westerham, Kent\\ 15th October.
\end{flushright}

\noindent Dear Sir,\\

\indent The disagreement subsisting between yourself and my late honoured father always gave me much uneasiness; and, since I have had the misfortune to lose him, I have frequently wished to heal the breach: but, for some time, I was kept back by my own doubts, fearing lest it might seem disrespectful to his memory for me to be on good terms with anyone with whom it had always pleased him to be at variance.\\
\end{quotation}

—»There, Mrs Bennet.«—

\begin{quotation}
My mind, however, is now made up on the subject; for, having received ordination at Easter, I have been so fortunate as to be distinguished by the patronage of the Right Honourable Lady Catherine de Bourgh, widow of Sir Lewis de Bourgh, whose bounty and beneficence has preferred me to the valuable rectory of this parish, where it shall be my earnest endeavour to demean myself with grateful respect towards her Ladyship, and be ever ready to perform those rites and ceremonies which are instituted by the Church of England. As a clergyman, moreover, I feel it my duty to promote and establish the blessing of peace in all families within the reach of my influence; and on these grounds I flatter myself that my present overtures of good-will are highly commendable, and that the circumstance of my being next in the entail of Longbourn estate will be kindly overlooked on your side, and not lead you to reject the offered olive branch. I cannot be otherwise than concerned at being the means of injuring your amiable daughters, and beg leave to apologize for it, as well as to assure you of my readiness to make them every possible amends; but of this hereafter. If you should have no objection to receive me into your house, I propose myself the satisfaction of waiting on you and your family, Monday, November 18th, by four o'clock, and shall probably trespass on your hospitality till the Saturday se'nnight following, which I can do without any inconvenience, as Lady Catherine is far from objecting to my occasional absence on a Sunday, provided that some other clergyman is engaged to do the duty of the day.\\

{\indent I remain, dear sir, with respectful compliments to your lady and daughters, your well-wisher and friend,}\\
\begin{flushright}
\textsc{William Collins.}
\end{flushright}
\end{quotation}


»At four o'clock, therefore, we may expect this peace-making gentleman,« said Mr Bennet, as he folded up the letter. »He seems to be a most conscientious and polite young man, upon my word; and, I doubt not, will prove a valuable acquaintance, especially if Lady Catherine should be so indulgent as to let him come to us again.«

»There is some sense in what he says about the girls, however; and, if he is disposed to make them any amends, I shall not be the person to discourage him.«

»Though it is difficult,« said Jane, »to guess in what way he can mean to make us the atonement he thinks our due, the wish is certainly to his credit.«

Elizabeth was chiefly struck with his extraordinary deference for Lady Catherine, and his kind intention of christening, marrying, and burying his parishioners whenever it were required.

»He must be an oddity, I think,« said she. »I cannot make him out. There is something very pompous in his style. And what can he mean by apologizing for being next in the entail? We cannot suppose he would help it, if he could. Can he be a sensible man, sir?«

»No, my dear; I think not. I have great hopes of finding him quite the reverse. There is a mixture of servility and self-importance in his letter which promises well. I am impatient to see him.«

»In point of composition,« said Mary, »his letter does not seem defective. The idea of the olive branch perhaps is not wholly new, yet I think it is well expressed.«

To Catherine and Lydia neither the letter nor its writer were in any degree interesting. It was next to impossible that their cousin should come in a scarlet coat, and it was now some weeks since they had received pleasure from the society of a man in any other colour. As for their mother, Mr Collins's letter had done away much of her ill-will, and she was preparing to see him with a degree of composure which astonished her husband and daughters.

Mr Collins was punctual to his time, and was received with great politeness by the whole family. Mr Bennet indeed said little; but the ladies were ready enough to talk, and Mr Collins seemed neither in need of encouragement, nor inclined to be silent himself. He was a tall, heavy-looking young man of five-and-twenty. His air was grave and stately, and his manners were very formal. He had not been long seated before he complimented Mrs Bennet on having so fine a family of daughters, said he had heard much of their beauty, but that, in this instance, fame had fallen short of the truth; and added, that he did not doubt her seeing them all in due time well disposed of in marriage. This gallantry was not much to the taste of some of his hearers; but Mrs Bennet, who quarrelled with no compliments, answered most readily,—

»You are very kind, sir, I am sure; and I wish with all my heart it may prove so; for else they will be destitute enough. Things are settled so oddly.«

»You allude, perhaps, to the entail of this estate.«

»Ah, sir, I do indeed. It is a grievous affair to my poor girls, you must confess. Not that I mean to find fault with \textit{you}, for such things, I know, are all chance in this world. There is no knowing how estates will go when once they come to be entailed.«

»I am very sensible, madam, of the hardship to my fair cousins, and could say much on the subject, but that I am cautious of appearing forward and precipitate. But I can assure the young ladies that I come prepared to admire them. At present I will not say more, but, perhaps, when we are better acquainted\longdash«

He was interrupted by a summons to dinner; and the girls smiled on each other. They were not the only objects of Mr Collins's admiration. The hall, the dining-room, and all its furniture, were examined and praised; and his commendation of everything would have touched Mrs Bennet's heart, but for the mortifying supposition of his viewing it all as his own future property. The dinner, too, in its turn, was highly admired; and he begged to know to which of his fair cousins the excellence of its cookery was owing. But here he was set right by Mrs Bennet, who assured him, with some asperity, that they were very well able to keep a good cook, and that her daughters had nothing to do in the kitchen. He begged pardon for having displeased her. In a softened tone she declared herself not at all offended; but he continued to apologize for about a quarter of an hour.